\documentclass[12pt,twoside]{report}
\usepackage[headheight=15pt]{geometry}
\usepackage{amsfonts, amsmath}
\usepackage[T1]{fontenc}
\usepackage{mathrsfs, enumitem}
\usepackage{hyperref}
\usepackage[utf8]{inputenc}
\usepackage{amssymb}
\usepackage{amsfonts}
\usepackage{amsmath}
\usepackage{amsthm}
\usepackage{color}
\usepackage{hyperref}
\usepackage{csquotes}
\usepackage{tikz-cd}
\usepackage{lipsum}
\usepackage{stmaryrd}
\usepackage{cancel}

\newtheorem{theorem}{Theorem}[subsection]
\newtheorem{lemma}[theorem]{Lemma}
\newtheorem{claim}[theorem]{Claim}
\newtheorem{proposition}[theorem]{Proposition}
\newtheorem{corollary}[theorem]{Corollary}
\newtheorem{fact}[theorem]{Fact}
\newtheorem{notation}[theorem]{Notation}
\newtheorem{observation}[theorem]{Observation}
\newtheorem{conjecture}[theorem]{Conjecture}
\newtheorem{exercise}[theorem]{Exercise}

\theoremstyle{definition}
\newtheorem{definition}[theorem]{Definition}
\newtheorem{example}[theorem]{Example}
\numberwithin{equation}{subsection}

\theoremstyle{remark}
\newtheorem{remark}[theorem]{Remark}
\theoremstyle{plain}
\newcommand{\ignore}[1]{}

% section symbol
%\renewcommand{\thesection}{\S\arabic{section}}

% \renewcommand{\Pr}{{\bf Pr}}
% \newcommand{\Prx}{\mathop{\bf Pr\/}}
% \newcommand{\E}{{\bf E}}
% \newcommand{\Ex}{\mathop{\bf E\/}}
% \newcommand{\Var}{{\bf Var}}
% \newcommand{\Varx}{\mathop{\bf Var\/}}
% \newcommand{\Cov}{{\bf Cov}}
% \newcommand{\Covx}{\mathop{\bf Cov\/}}

% shortcuts for symbol names that are too long to type
\newcommand{\eps}{\epsilon}
\newcommand{\lam}{\lambda}
\renewcommand{\l}{\ell}
\newcommand{\la}{\langle}
\newcommand{\ra}{\rangle}
\newcommand{\wh}{\widehat}
\newcommand{\wt}{\widetilde}

% % "blackboard-fonted" letters for the reals, naturals etc.
\newcommand{\R}{\mathbb R}
\newcommand{\N}{\mathbb N}
\newcommand{\Z}{\mathbb Z}
\newcommand{\F}{\mathbb F}
\newcommand{\Q}{\mathbb Q}
\newcommand{\C}{\mathbb C}
\newcommand{\D}{\mathbb D}
\newcommand{\T}{\mathbb T}

% % operators that should be typeset in Roman font
% \newcommand{\poly}{\mathrm{poly}}
% \newcommand{\polylog}{\mathrm{polylog}}
% \newcommand{\sgn}{\mathrm{sgn}}
% \newcommand{\avg}{\mathop{\mathrm{avg}}}
% \newcommand{\val}{{\mathrm{val}}}

% % complexity classes
% \renewcommand{\P}{\mathrm{P}}
% \newcommand{\NP}{\mathrm{NP}}
% \newcommand{\BPP}{\mathrm{BPP}}
% \newcommand{\DTIME}{\mathrm{DTIME}}
% \newcommand{\ZPTIME}{\mathrm{ZPTIME}}
% \newcommand{\BPTIME}{\mathrm{BPTIME}}
% \newcommand{\NTIME}{\mathrm{NTIME}}

% values associated to optimization algorithm instances
\newcommand{\Opt}{{\mathsf{Opt}}}
\newcommand{\Alg}{{\mathsf{Alg}}}
\newcommand{\Lp}{{\mathsf{Lp}}}
\newcommand{\Sdp}{{\mathsf{Sdp}}}
\newcommand{\Exp}{{\mathsf{Exp}}}

% if you think the sum and product signs are too big in your math mode; x convention
% as in the probability operators
\newcommand{\littlesum}{{\textstyle \sum}}
\newcommand{\littlesumx}{\mathop{{\textstyle \sum}}}
\newcommand{\littleprod}{{\textstyle \prod}}
\newcommand{\littleprodx}{\mathop{{\textstyle \prod}}}

% horizontal line across the page
\newcommand{\horz}{
\vspace{-.4in}
\begin{center}
\begin{tabular}{p{\textwidth}}\\
\hline
\end{tabular}
\end{center}
}

% calligraphic letters
\newcommand{\calA}{{\cal A}}
\newcommand{\calB}{{\cal B}}
\newcommand{\calC}{{\cal C}}
\newcommand{\calD}{{\cal D}}
\newcommand{\calE}{{\cal E}}
\newcommand{\calF}{{\cal F}}
\newcommand{\calG}{{\cal G}}
\newcommand{\calH}{{\cal H}}
\newcommand{\calI}{{\cal I}}
\newcommand{\calJ}{{\cal J}}
\newcommand{\calK}{{\cal K}}
\newcommand{\calL}{{\cal L}}
\newcommand{\calM}{{\cal M}}
\newcommand{\calN}{{\cal N}}
\newcommand{\calO}{{\cal O}}
\newcommand{\calP}{{\cal P}}
\newcommand{\calQ}{{\cal Q}}
\newcommand{\calR}{{\cal R}}
\newcommand{\calS}{{\cal S}}
\newcommand{\calT}{{\cal T}}
\newcommand{\calU}{{\cal U}}
\newcommand{\calV}{{\cal V}}
\newcommand{\calW}{{\cal W}}
\newcommand{\calX}{{\cal X}}
\newcommand{\calY}{{\cal Y}}
\newcommand{\calZ}{{\cal Z}}

% bold letters (useful for random variables)
%----------------------------------------------
% \renewcommand{\a}{{\boldsymbol a}}
% \renewcommand{\b}{{\boldsymbol b}}
% \renewcommand{\c}{{\boldsymbol c}}
% \renewcommand{\d}{{\boldsymbol d}}
% \newcommand{\e}{{\boldsymbol e}}
% \newcommand{\f}{{\boldsymbol f}}
% \newcommand{\g}{{\boldsymbol g}}
% \newcommand{\h}{{\boldsymbol h}}
% \renewcommand{\i}{{\boldsymbol i}}
% \renewcommand{\j}{{\boldsymbol j}}
% \renewcommand{\k}{{\boldsymbol k}}
% \newcommand{\m}{{\boldsymbol m}}
% \newcommand{\n}{{\boldsymbol n}}
% \renewcommand{\o}{{\boldsymbol o}}
% \newcommand{\p}{{\boldsymbol p}}
% \newcommand{\q}{{\boldsymbol q}}
% \renewcommand{\r}{{\boldsymbol r}}
% \newcommand{\s}{{\boldsymbol s}}
% \renewcommand{\t}{{\boldsymbol t}}
% \renewcommand{\u}{{\boldsymbol u}}
% \renewcommand{\v}{{\boldsymbol v}}
% \newcommand{\w}{{\boldsymbol w}}
% \newcommand{\x}{{\boldsymbol x}}
% \newcommand{\y}{{\boldsymbol y}}
% \newcommand{\z}{{\boldsymbol z}}
% \newcommand{\A}{{\boldsymbol A}}
% \newcommand{\B}{{\boldsymbol B}}
% \newcommand{\C}{{\boldsymbol C}}
% \newcommand{\D}{{\boldsymbol D}}
% \newcommand{\E}{{\boldsymbol E}}
% \newcommand{\F}{{\boldsymbol F}}
% \newcommand{\G}{{\boldsymbol G}}
% \renewcommand{\H}{{\boldsymbol H}}
% \newcommand{\I}{{\boldsymbol I}}
% \newcommand{\J}{{\boldsymbol J}}
% \newcommand{\K}{{\boldsymbol K}}
% \renewcommand{\L}{{\boldsymbol L}}
% \newcommand{\M}{{\boldsymbol M}}
% \renewcommand{\O}{{\boldsymbol O}}
% \renewcommand{\P}{{\mathbb{P}}}
% \newcommand{\Q}{{\boldsymbol Q}}
% \newcommand{\R}{{\boldsymbol R}}
% \renewcommand{\S}{{\boldsymbol S}}
% \newcommand{\T}{{\boldsymbol T}}
% \newcommand{\U}{{\boldsymbol U}}
% \newcommand{\V}{{\boldsymbol V}}
% \newcommand{\W}{{\boldsymbol W}}
% \newcommand{\X}{{\boldsymbol X}}
% \newcommand{\Y}{{\boldsymbol Y}}
% \newcommand{\Z}{{\boldsymbol Z}}

% script letters
\newcommand{\scrA}{{\mathscr A}}
\newcommand{\scrB}{{\mathscr B}}
\newcommand{\scrC}{{\mathscr C}}
\newcommand{\scrD}{{\mathscr D}}
\newcommand{\scrE}{{\mathscr E}}
\newcommand{\scrF}{{\mathscr F}}
\newcommand{\scrG}{{\mathscr G}}
\newcommand{\scrH}{{\mathscr H}}
\newcommand{\scrI}{{\mathscr I}}
\newcommand{\scrJ}{{\mathscr J}}
\newcommand{\scrK}{{\mathscr K}}
\newcommand{\scrL}{{\mathscr L}}
\newcommand{\scrM}{{\mathscr M}}
\newcommand{\scrN}{{\mathscr N}}
\newcommand{\scrO}{{\mathscr O}}
\newcommand{\scrP}{{\mathscr P}}
\newcommand{\scrQ}{{\mathscr Q}}
\newcommand{\scrR}{{\mathscr R}}
\newcommand{\scrS}{{\mathscr S}}
\newcommand{\scrT}{{\mathscr T}}
\newcommand{\scrU}{{\mathscr U}}
\newcommand{\scrV}{{\mathscr V}}
\newcommand{\scrW}{{\mathscr W}}
\newcommand{\scrX}{{\mathscr X}}
\newcommand{\scrY}{{\mathscr Y}}
\newcommand{\scrZ}{{\mathscr Z}}

\newcommand{\im}{{\text{im }}}
\newcommand{\ip}[1]{\left\langle #1 \right\rangle}
\newcommand{\norm}[1]{\left\lVert #1 \right\rVert}
\newcommand{\abs}[1]{\left\lvert #1 \right\rvert}

\usepackage{fancyhdr}

\fancypagestyle{myheader}{
    \fancyhead{} %clear all headers
    \fancyfoot{} %clear all footers
    \fancyhead[CO]{\textsc{\leftmark}} %How to insert the Chapter names like "MyChapter"?
    \fancyfoot[C]{\thepage}
    \fancyhead[CE]{\textsc{Ashish Kujur}}
    \renewcommand{\headrulewidth}{0.4pt}
    \renewcommand{\footrulewidth}{0.0pt}
}
\pagestyle{myheader}
\renewcommand{\chaptermark}[1]{\markboth{\chaptername\ \thechapter:\ #1}{}}

\title{
{Study of Closed Ideals of Disc Algebra}\\
{\large Indian Institute of Science Education and Research -- Thiruvananthapuram}\\ \vspace{1cm}
{\includegraphics[height=0.1\textheight]{iiser_logo.png}}
}
\author{Ashish Kujur}
\date{Last Updated: \today}


\begin{document}
\maketitle \tableofcontents
\chapter{Analytic and Harmonic Functions}
Let us begin by fixing the notation for the rest of the report. The \textit{open unit disc} will be denoted by $\D$ and is defined to the subset of the complex plane of those elements whose absolute value is less than $1$. That is,
\begin{equation*}
    \D = \left\{ z \in \C : \abs{z} < 1 \right\}.
\end{equation*}
The boundary of the open unit disc, called, the \textit{unit circle} is denotes by $\T$. That is, 
\begin{equation*}
    \T = \left\{ z \in \C : \abs{z} = 1 \right\}.
\end{equation*}

We will assume that the the unit circle has the normalised Lebesgue measure. The set of all complex measures denoted by $\calM \left( \T \right)$ is a Banach space with total variation norm. For a proof of this fact, we advise the reader to refer \cite{cohn2013measure}, \cite{axler2020measure} and \cite{rudin1987real}.

As usual, $\Omega$ will denote an open connected subset of $\C$. We will say that a function $U : \Omega \to \C$ is \textit{harmonic} on $\Omega$ if it satisfies the Laplace equation
\begin{equation*}
    \triangledown ^{2} U = \frac{\partial ^{2} U}{\partial x^{2}} + \frac{\partial ^{2} U}{\partial y^{2}} =0
\end{equation*}
for each point in $\Omega$.


Before we proceed, recall a basic fact from complex analysis which states that a real valued function on $\D$ is harmonic iff it is the real part an analytic function.

\section{The Poisson Kernel}
The famous Dirchlet problem asks the following:


Given $G \subset \R ^{2}$ and $f: \partial G \to \C$ and $f: \partial G \to \C$. Find a continuous function $u: \overline{G} \to \C$ such that $u\mid _{G}$  is harmonic and $u\mid_{\partial G} =f$.

To solve the above problem, we make the following definition:

\begin{definition}
    For each $0\le r < 1$, we define $P_r : \T \to \C$ by
    \begin{equation*}
	P_{r} \left( e^{it} \right) = \frac{1-r^{2}}{1+r^{2}-2r\cos t}
    \end{equation*}
    for each $e^{it} \in \T$.
    The family of functions $\left\{ P_{r} \right\}_{r\in [0,1)}$ is called the \textit{Poisson kernel} on the open unit disc $\D$.
    \label{def:Poisson-kernel}
\end{definition}

It can be shown that that for each $r<1$, we have that 
\begin{equation*}
    P_{r}\left( e^{it} \right) = \sum_{n=-\infty}^{\infty} r^{\abs{n}}e^{int}
\end{equation*}
holds for each $e^{it} \in \T$.

We will soon see that if $u \in \calC \left( \T \right)$, that is, continuous on $\T$, the so called "Poisson integral" of $u$ solves the Dirichlet problem. We postpone this work to the next section. However the following proposition may lead the reader to believe why this may be true:
\begin{proposition}
    Let $u: \overline{\D} \to \C$ be a harmonic function, this means, that $u$ is harmonic in an open subset containing $\overline{D}$. Then we have that
    \begin{equation}
	u\left( re^{i\theta} \right) = \frac{1}{2\pi} \int_{-\pi}^{\pi} u\left( e^{it} \right) P_{r} \left( e^{i\left( \theta -t \right)} \right) dt
    \end{equation}
    \label{prop:poisson-integral-formula-for-harmonic-function}
\end{proposition}

Albeit the proof of this proposition is not too hard, we ask the reader to seek a proof of this fact in \cite{hoffman2007banach}.

\section{Boundary Values}

The Poisson kernel is an "approximate identity". We just deal with the approximate identities on the circle:

\begin{definition}
    Let $\varphi_{i}$ be a subset of $L^{1} \left( \T \right)$ where $i$ varies over a directed set. We say that $\left\{ \varphi_{i} \right\}$ is an approximate identity on $\T$ if the following are satisfied:
    \begin{enumerate}
	\item for all $i$,
	    \begin{equation*}
		\frac{1}{2\pi} \int_{-\pi}^{\pi} \varphi_{i} \left( e^{it} \right) dt =1
	    \end{equation*}
	\item
	    \begin{equation*}
		C_{\varphi} = \sup_{i} \left( \frac{1}{2\pi} \int_{-\pi}^{\pi} \abs{\varphi_{i} \left( e^{it} \right)dt} \right) \le \infty
	    \end{equation*}
	\item for each fixed $\delta$, $0<\delta<\pi$, 
	    \begin{equation*}
		\lim_{i} \int_{\delta \le \abs{t} \le \pi} \abs{\varphi_{i} \left( e^{it} \right)} dt = 0.
	    \end{equation*}
    \end{enumerate}
    A family of approximate identity on $\T$ is called \textit{positive approximate identity} if all the function are positively valued, that is, 
    \begin{equation*}
	\varphi_{i}\left( e^{it} \right) \ge 0
    \end{equation*}
    for each $i$.
\end{definition}
As mentioned, the Poisson kernel is an example of approximate identity.The other well known approximate identity is Fejer's kernel on $\T$ which is essential in theory of convergence of Fourier series on $\T$. We refer the reader to \cite{katznelson2004introduction}.


The aforementioned theorems are esssential to the studying the interaction of $\D$ and the $\T$. We ask the reader to see \cite{mashreghi2009representation} for proofs for the sake of brevity. Although these theorems are done more generally, we really only care about the special case of the Poisson kernel.

\begin{theorem}[Weak* convergence of measures]
    Let $\left\{ \varphi _{i} \right\}_{i}$ be an approximate identity on $\T$ and let $\mu \in \calM \left( \T \right)$. Then for all $i$, $\varphi_{i} * \mu  \in L^{1} \left( \T \right)$ with
    \begin{equation*}
	\norm{\varphi _{i} * \mu}_{1} \le C_{\varphi} \lVert \mu \rVert
    \end{equation*}
    and
    \begin{equation*}
	\norm{\mu} \le \sup_{i} \norm{\varphi_{i} * \mu}_{1}\text{.}
    \end{equation*}
    Moreover, the measures $d\mu_{i} = \left( \varphi_{i} * \mu \right) \left( e^{it} \right) dt/2\pi$ converge to $d\mu \left( e^{it} \right)$ in the weak* topology, i.e.
    \begin{equation*}
	\lim_{i} \frac{1}{2\pi} \int_{-\pi}^{\pi} f\left( e^{it} \right) \left( \varphi_{i} * \mu \right) \left( e^{it} \right) dt = \int_{\T} \varphi\left( e^{it} \right) d\mu \left( e^{it} \right)
    \end{equation*}
    for all $f \in \calC \left( \T \right)$.
    \label{thm:weak-star-measures}
\end{theorem}
\begin{theorem}[Convergence in norm]
    Let $\left\{ \varphi _{i} \right\}_{i}$ be an approximate identity on $\T$ and let $f \in L^{p} \left( \T \right)$ with $p \in [1, \infty)$. Then for all $i$, $\varphi_{i} * f  \in L^{p} \left( \T \right)$ with
    \begin{equation*}
	\norm{\varphi _{i} * f}_{p} \le C_{\varphi} \lVert f \rVert _{p}
    \end{equation*}
    and
    \begin{equation*}
	\lim_{i} \norm{\varphi_{i} * f -f}_{p} = 0 \text{.}
	    \end{equation*} 
    \label{thm:convergence-in-Lp}
\end{theorem}


\begin{theorem}[Weak* convergence of bounded functions]
    Let $\left\{ \varphi _{i} \right\}_{i}$ be an approximate identity on $\T$ and let $f\in L^{\infty} \left( \T \right)$. Then for all $i$, $\varphi_{i} * \mu  \in \calC \left( \T \right)$ with
    \begin{equation*}
	\norm{\varphi _{i} * \mu}_{\infty} \le C_{\varphi} \lVert \mu \rVert _{\infty}
    \end{equation*}
    and
    \begin{equation*}
	\norm{f}_{+\infty} \le \sup_{i} \norm{\varphi_{i} * f}_{\infty}\text{.}
    \end{equation*}
    Moreover, $\varphi_{i} * f$ converge to $f$ in the weak* topology, i.e.
    \begin{equation*}
    \lim_{i} \int_{-\pi}^{\pi} g\left( e^{it} \right) \left( \varphi_{i} * f \right) \left( e^{it} \right) dt = \int_{\T} g\left( e^{it} \right) f \left( e^{it} \right) dt 
    \end{equation*}
    for all $g \in L^{1} \left( \T \right)$.
    \label{thm:weak-star-infinity}
\end{theorem}

\subsection{The specific case of Poisson Integral}

\begin{definition}[Poisson integral of some function or measure]
    Let $\tilde{f} : \D \to \C$ be a harmonic function. Then $\tilde{f}$ is said to be the \textit{Poisson integral} of the function $f : \T \to \C$ if
    \begin{equation*}
	\tilde{f} (re^{i\theta}) = \frac{1}{2\pi} \int_{T} f\left( e^{it} \right) P_{r} \left( e^{i\left( \theta-t \right)} \right) dt
    \end{equation*}
    In such a case, we will denote the function $\tilde{f}$ by $P[f]$.
    Similarly, $f$ is said to be the \textit{Poisson integral} of a complex measure $\mu$ on $T$ if
\begin{equation*}
    \tilde{f} (re^{i\theta}) = \frac{1}{2\pi} \int_{T} P_{r} \left( e^{i\left( \theta-t \right)} \right) d\mu\left( e^{it} \right)
    \end{equation*}In such a case, we will denote the function $\tilde{f}$ by $P[\mu]$.
    \label{def:Poisson-Integral-Of-Some-Function-Or-Measure}
\end{definition}

This follows immediately by noticing that the family of Poisson kernel is indeed an approximate identity. 
\begin{theorem}
    Let $f : \D \to \C$ be a harmonic function. Define for each $r\in [0,1)$, the function $f_{r} : \T \to \C$ by
    \begin{equation*}
	f_{r} \left( e^{i\theta} \right) = f\left( re^{i\theta} \right)
    \end{equation*}
    The following statements holds:
    \begin{enumerate}
	\item If $1 < p \le \infty$ then $f=P[g]$ for some $g \in L^{p} [g]$ iff for each $r > 0$, $\norm{f_{r}}_{p} < +\infty$ .
	\item If p=1 then $f=P[g]$ for some $g \in L^{p} [g]$ iff $f_{r}$ converge in the $L^{1}$ norm.
	\item $f=P[\mu]$ for some $\mu \in \calM (\T)$ iff for each $r > 0$, $\norm{f_{r}}_{1} < +\infty$ 
    \end{enumerate}
    \label{thm:convergence-Poisson}
\end{theorem}
\section{Fatou's Theorem}

The next theorem requires the reader to know what derivative of measures are. The reader can get review of these concepts in \cite{rudin1987real} and \cite{cohn2013measure}.

\begin{theorem}[Fatou (1905)]
    Let $\mu$ be a complex measure on the unit circle $\T$, and let $f: \D \to \C$ be the harmonic function defined by
    \begin{align*}
	f(re^{i\theta}) = \frac{1}{2\pi} \int_{\T} P_{r} \left( e^{i\left( \theta-t \right)} \right) d\mu \left( e^{it} \right)
    \end{align*}

    Let $e^{i\theta_{0}}$ be any point where $\mu$ is differentiable with respect to the normalised Lebesgue measure. Then
    \begin{equation*}
	\lim_{r\to 1} f\left( re^{i\theta_{0}} \right) = \left( \frac{d\mu}{d\theta} \right) \left( e^{i\theta _{0}} \right) = \mu ' \left( e^{i\theta _{0}} \right)
    \end{equation*}
    In fact, $f(re^{i\theta}) \to \mu ' \left( e^{i\theta_{0}} \right)$ as $re^{i\theta}$ approaches $e^{i\theta_{0}}$ along any path in the open disc within the region of the form $\abs{\theta - \theta_{0}} \le c \left( 1-r \right)$ for some $c> 0$. 
    \label{thm:Fatou-1906}
\end{theorem}

A proof of the above theorem is quite technical and can be found in \cite{koosis1998introduction}, \cite{hoffman2007banach} and \cite{axler2013harmonic}.

\begin{corollary}
    Let $\mu$ be a complex measure on $\T$. Then $P[\mu]$ has nontangential limits equal everywhere to the Radon Nikodym derivative of $\mu$ with respect to the normalised Lebesgue measure.
\end{corollary}
\begin{proof}
    Let $\mu$ be a complex measure on the unit circle. Then by LebesgueDecomposition Theorem, we have that

    \begin{equation*}
	d\mu = \frac{1}{f} dt + d\mu_{s}
    \end{equation*}
    where $d\mu_{s}$ is a singular measure. 
    Since $\mu$ is a complex measure, hence, its real and imaginary parts are finite measures, we have that $\mu$ is differentiable almost everywhere and we have that
    \begin{equation*}
	\frac{d\mu}{dt} = \frac{1}{2\pi}f
    \end{equation*}
    almost everywhere.

    Applying Fatou's theorem \ref{thm:Fatou-1906}, this follows immediately.
\end{proof}

\begin{corollary}
    Let $f : \T \to \C$ be $L^{1}$. Then $P[f]$ has nontangential limits at almost everywhere and these limits equal to $f$ almost everywhere.
    \label{cor:L1-implies-Poisson-limits}
\end{corollary}
\begin{proof}
    This corollary follows immediately by considering the measure on the circle given by
    \begin{equation*}
	d\mu = f \frac{dt}{2\pi}
    \end{equation*}
    Applying the previous corollary, we are done.
\end{proof}

The next corollary tells us that if the $p$ norm of $f_{r}$ is uniform bounded as $r$ varies, the $f$ is a Poisson integral of a $L^{p}$ function on the circle or a complex measure on the circle depending upon what $p$ is.
\begin{corollary}
    Let $f: \D \to \C$ be a harmonic function and $1\le p <\infty$. Suppose that for all $0\le r < 1$, we have that
    \begin{equation*}
	\norm{f_{r}}_{p} < M< +\infty
    \end{equation*}
    Then for almost every $\theta$ the radial limits 
    \begin{equation*}
	\tilde {f} (e^{i\theta} ) = \lim_{r\to 1} f\left( re^{i\theta} \right)
    \end{equation*}
    exist and define a function $\tilde f$ in $L^{p} \left( \T \right)$. The following also holds:
    \begin{enumerate}
	\item If $p>1$ then $f=P[\tilde{f}]$.
	\item If $p=1$ then $f=P[\mu]$ for some complex measure $\mu$ whose absolutely continuous part is $fd\theta$.
	\item IF $f$ is bounded then the boundary values exist almost everywhere and define a bounded measurable function $\tilde{f}$ on $\T$ such that $f=P[\tilde{f}]$.
    \end{enumerate}
    \label{cor:imp-Fatou}
\end{corollary}



\begin{proof}
    Suppose that for each $r\in [0,1)$, we have $\norm{f_{r}}_{p} < +\infty$. We need to prove that for almost every $\theta$, $\lim_{r\to 1} f\left( re^{i\theta} \right)$ exists. Then by Theorem \ref{thm:convergence-Poisson}, we have that $f=P[g]$ for some $g\in L^{p} \left( \T \right)$. Since $L^{p} \left( \T \right) \subset L^{1} \left( \T \right)$, we can use the previous corollary. By the previous corollary, we have that $P[g]$ has nontangential limits almost everywhere, we have that
    \begin{equation}
	\tilde{f} \left( e^{i\theta} \right) = \lim_{r\to 1} f(re^{i\theta}) = \lim_{r\to 1} P[g] \left( re^{i\theta} \right)
	\label{eqn:radial-limits}
    \end{equation}
    exists almost everywhere.

    Now we proceed to prove part $(1)$. Also by Theorem \ref{thm:convergence-Poisson}, we have that $f=P[g]$ for some $g\in L^{p} \left( \T \right)$. Hence, we have that by Equation \ref{eqn:radial-limits} that $\tilde{f} (e^{i\theta}) =  \lim_{r\to 1} P[g] \left( re^{i\theta} \right)$ holds at almost every $\theta$.

    Also, by the previous corollary, $\lim_{r\to 1} P[g] \left( re^{i\theta} \right) = g(e^{i\theta})$ for almost every $\theta$. Hence, we have that $\tilde{f} = g$.

\end{proof}

\begin{corollary}
    Let $f: \D \to \R_{\ge 0}$ be a harmonic function. Then $f$ has nontangential limits at almost every point of $\T$.
    \label{cor:nontangential-limits}
\end{corollary}
\begin{proof}
    The proof just relies on \ref{prop:poisson-integral-formula-for-harmonic-function}.
\end{proof}

Let $h \left( \D \right)$ denote the set of all harmonic functions on $\D$. Let $p\in [1,\infty]$. Define
\begin{equation*}
    h^{p} \left( \D \right) = \left\{ f\in h \left( \D \right) \, \mid \, \left\{ f_{r} \right\}_{0\le r < 1} \text{ is uniformly bounded in } L^{p} \text{ norm }\right\}
\end{equation*}
We define a norm on $h^{p} \left( \D \right)$ by
\begin{equation*}
    \norm{f}_{h^{p} \left( \D \right)} = \sup_{0\le r < 1} \norm{f_{r}}_{L^p \left( \D \right)} =\begin{cases} \sup_{0\le r < 1} \left( \frac{1}{2\pi} \int_{-\pi}^{\pi} \abs{f\left( re^{i\theta}\right)}^{p}d\theta   \right)^{\frac{1}{p}} & \text{if } p \in [1, \infty) \\
	\sup_{0\le r < 1} \norm{f(re^{i\theta})}_{L^{\infty} \left( \D \right)} & \text{ if } p=\infty
    \end{cases}
\end{equation*}

It is easy to see why $\norm{f} < +\infty$ for any $f\in h^{p} \left( D \right)$. So we now proceed to show that $h^{p} \left( D \right)$ is a Banach space with this norm. First we show that it is indeed a normed linear space.

Clearly, $h \left( \D \right)$ is a vector space. To show that $h^{p} \left( \D \right)$ is a vector space, it suffices to check that $h^{p} \left( \D \right)$ is a subspace.

Let $f,g \in h^{p} \left( \D \right)$ and let $\alpha \in \C$. Then for any $r\in [0,1)$, we have that 
\begin{align*}
    \norm{(f+\alpha g)_{r}}_{p} &= \norm{f_{r} + \alpha g_{r}} \\
    &= \norm{f_{r}}_{p} + \alpha \norm{g_{r}}_{p}
\end{align*}
Take note of the use of Holder's inequality. After this is done, since $\left\{ f_{r} \right\}_{r\in [0,1)}$ and $\left\{ g_{r} \right\}_{r\in [0,1)}$ is uniformly bounded, we have that $\left\{ f+ \alpha g \right\}_{r\in [0,1)}$ is uniformly bounded in $L^{p}$ norm.

Now, we need to show that it is a normed linear space but this follows almost immediately.

To show that it is a Banach space, we show that

\begin{theorem}
    Let $p\in [1, \infty]$. If $u\in L^{p} \left( \T \right)$ then $f=P * u \in h^{p} \left( \D \right)$ and $\norm{f}_{p} =\norm{u}_{p}$. If $\mu \in \calM \left( \T \right)$ then $f= P * \mu \in h^{1} \left( \D \right)$ and $\norm{f}_{1} = \norm{\mu}$.
    \label{thm:lp-and-hp}
\end{theorem}
\begin{proof}
    We consider the case $p \in [1, \infty )$. The other cases can be dealt similarly. Consider the map 
    \begin{equation*}
	u \stackrel{T}{\mapsto} U
    \end{equation*}
    where $U\left( re^{i\theta} \right) = \frac{1}{2\pi} \int_{-\pi}^{\pi} P_{r}\left( e^{i(\theta-t)} \right) u \left( e^{it} \right) dt$. By Theorem \ref{thm:convergence-in-Lp}, we have that $\norm{U} = \norm{u}_{p} < +\infty$. Hence $U \in h^{p} \left( \D \right)$.

    Linearity is obvious. We need to check injectivity and surjectivity.

    To check injectivity, let $u \in L^{p} \left( \T \right)$ and suppose that $T(u)=P[u]=0$. Now $\lim_{r\to 1} P[u] \left( re^{i\theta} \right) = u$ for almost $\theta$ by Corollary \ref{cor:L1-implies-Poisson-limits} and hence $u=0$ almost everywhere.

    Surjectivity is clear from Theorem \ref{thm:convergence-Poisson}.
    \end{proof}

%%%%%%%%%%%%%%%%%%%%%%%%%%%%%%%%%%%%%%%%%%%%%%%%%%%%%%%%%%%%%%%%%%%%%%%%%%%%%%

\section{Hardy Spaces -- \texorpdfstring{$H^p$}{\text{Hp}} spaces}

Let us denote the set of all analytic functions on $\D$ by $H \left( \D \right)$. Hence, $H \left( \D \right) \subset h \left( \D \right)$. For $p\in (0,\infty]$, we consider the \textit{Hardy classes} of analytic functions on the unit disc

\begin{equation*}
    H^{p} \left( \D \right) = \left\{ F \in H \left( \D \right) \, \mid \, \norm{F} _{p} < \infty \right\}
\end{equation*}

Clearly, 
\begin{align*}
    H^{p} \left( \D \right) \subset h ^{p} \left( \D \right)
\end{align*}

We will see that $H ^{p} \left( \D \right)$, $1\le p \le +\infty$, is also a Banach spaces isomorphic to a closed subspace of $L^{p} \left( \T \right)$ denotes by $H^{p} \left( \T \right)$.

To prove that $H^{p} \left( \D \right)$ is a closed subspace of $h^{p} \left( \D \right)$, we are going to identify $H^{p} \left( \D \right)$ with the closed subspace
\begin{align*}
    \left\{ u\in L^{p} \left( \T \right) : \int_{-\pi}^{\pi} u\left( e^{it} \right) e^{ikt} = 0 \textit{ for all } k \in \N \right\}
\end{align*}

Let $\left\{ u_{n} \right\}$ be a sequence of functions in the above subspace; suppose that $\left\{ u_{n} \right\}$ converge to $u\in L^{p} \left( \T \right)$. Now, let $k\in \N$ be arbitrary. Since $\left\{ u_{n} \right\}$ converge to $u$ in $p$-norm, we have that $\left\{ u_{n} \right\}$ converge to $u$ in $1$-norm. Hence we have the following:
\begin{align*}
    \abs{\int_{-\pi}^{\pi} u_{n} \left( e^{it} \right) e^{ikt} dt - \int_{-\pi}^{\pi} u \left( e^{it} \right) e^{ikt} dt } &\le \int_{-\pi}^{\pi} \abs{u_{n} \left( e^{it} \right) - u\left( e^{it} \right)} dt
\end{align*}
From the above inequality, it is evident that $u$ is in the subspace mentioned above.

For the above reasons, we will not distinguish between the $H^{p} \left( \D \right)$ and $H^{p}\left( \T \right)$ and simply write $H^{p}\left( \D \right)$. The case $p=1$ needs a special mention. We have not dealt with it yet. We saw that if $f \in H^{1}\left( \D \right)$ then we have that 
\begin{equation*}
    f = P[\mu]
\end{equation*}
for some complex measure $\mu$. We do not recover a $f\in L^{1} \left( \T \right)$ like the other case. In this case, we know that this measure $\mu$ satisfies the property that
\begin{equation*}
    \int_{-\pi}^{\pi} e^{int} dt = 0
\end{equation*}
for each $n=1,2,3,\ldots$. In other words, we can say the negative Fourier coefficient of $\mu$ is zero. We will call a such a complex measure to be \textit{analytic}.

In the next chapter, we will see that a theorem due to F and M Riesz says that complex analytic measure is absolutely continuous with respect to the normalised Lebesgue measure. That is, there is a $\tilde{f}\in L^{1} \left( \T \right)$ such that 
\begin{equation*}
    d\mu = \tilde{f} \frac{dt}{2\pi}.
\end{equation*}

In the proof of Theorem \ref{thm:lp-and-hp}, we can send $f$ to $\tilde{f}$ to establish a bijection and we will be done.

\subsubsection{Series Representation of Harmonic Functions}

The proof of the aforementioned theorem is used later and not necessary for our discussion. We refer the reader to \cite{mashreghi2009representation} for a proof.

\begin{theorem}
    Let $U$ be a harmonic on the disc $D_{R} = \left\{ |z| < R \right\}$. Then, for each $n\in \Z$, the quantity
    \begin{equation}
	a_{n} = \frac{\rho^{-\abs{n}}}{2\pi} \int_{-\pi}^{\pi} U\left( \rho e^{it} \right) e^{-int} dt \qquad \left( 0 < \rho < R \right)
	\label{eqn:srhf-1}
    \end{equation}
    is independent of $\rho$ and we have
    \begin{equation}
	U\left( re^{i\theta} \right) = \sum_{n=-\infty}^{\infty} a_{n} r^{\abs{n}} e^{in\theta} \qquad \left( re^{i\theta} \in \D \right)
	\label{eqn:srhf-2}
    \end{equation}
    The function
    \begin{equation}
	V\left( re^{i\theta} \right) = \sum_{n=-\infty}^{\infty} -i \text{sgn } (n) a_{n} r^{\abs{n}} e^{in\theta} \qquad \left( re^{i\theta} \in \D \right)
	\label{eqn:srhf-3}
    \end{equation}
    is the unique harmonic conjugate of $U$ such that $V\left( 0 \right)  = 0$. The series in \ref{eqn:srhf-2} and \ref{eqn:srhf-3} are absolutely and uniformly convergent on compact subsets of $D_{R}$
    \label{thm:srhf}
\end{theorem}

\chapter{The space \texorpdfstring{$H^1$}{\text{H1}}}
We use the following theorem several times in this chapter. We mention it here again for reference:
\begin{theorem}
    Let $u : \overline{\D} \to \C$ be a harmonic function. Then we have that
    \begin{equation*}
	u\left( re^{i\theta} \right) = \frac{1}{2\pi} \int_{-\pi}^{\pi} u\left( e^{it} \right) P_{r} \left( e^{i\left( \theta-t \right)} \right)
    \end{equation*}
    \label{thm:Poisson-Integral-Formula}
\end{theorem}

With this, we proceed to prove the F and M Riesz theorem and the Szego's theorem. The former helps us establish an isometry between $H^{1}\left( \D \right)$ and latter talks about the integrability of
\begin{equation*}
    \log |F|
\end{equation*}
where $|F|$ is a $H^{1}$ function on the $\T$. The integrability of $\log \abs{F}$ will help us characterise all "outer functions" but more on that in the later chapter. Like the $L^{p}$ spaces, $H^{1}$ contains all the other $H^{p}$ spaces for $p>1$ and in the final section of this chapter, we study more properties of $H^{1}$ functions.

\section{The Helson-Lowdenslager Approach}
Let $\calC \left( \overline {\D} \right)$ be the set of all continuous functions on $\overline{\D}$ and let $H\left( \D \right)$ be the set of all holomorphic functions on the open disc $\D$. We define $\calA =\calC \left( \overline{ \D } \right) \cap H \left( \D \right)$.

We show that $\calA$ is an uniformly closed algebra of $\calC \left( \overline{ \D } \right)$. Let $\left\{ f_{n} \right\}$ be a sequence in $\calA$ converging uniformly to $f\in \calC \left( \overline \D \right)$.

We recall Morera's Theorem for analytic functions at this point:
\begin{theorem}[Morera]
    A continuous, complex valued function $f : D \to \C$ that satisfies $\oint _{\gamma} f \left( z \right) dz = 0$ for any closed piecewise $C^{1}$ path $\gamma$ in $D$ must be holomorphic on $D$.
    \label{thm:morera-analytic}
\end{theorem}

A proof of Morera's theorem can be found in any text in Complex Analysis. For instance the reader may refer \cite{ash2014complex} or \cite{rudin1987real} for a proof whose statements of which are stated in more or less in a similar fashion.

We use the Morera theorem to prove what we want to prove. Now, let $C$ be any closed curve in $\D$. Then for any $n\in \mathbb N$,
\begin{align*}
\oint_{C}f_{n}  \left( z \right) dz = 0
\end{align*}
So, 
\begin{align*}
    \oint_{C} f(z) dz = \oint_{C} \lim_{n\to \infty} f_{n} \left( z \right) dz = \lim_{n\to \infty} \oint_{C} f_{n} \left( z \right) dz =0
\end{align*}
Since $C$ was arbitrary, $f$ must be holomorphic. This shows that $\calA$ is uniformly closed. The fact that it is an algebra is easy to check.

Now, note that since $\D$ is a compact metric space, we have that $\calC \left( \D \right)$ is a complete metric space with supremum metric. Since the supremum metric can also be induced by a norm, namely the supremum norm, we have that $\calC\left( \D \right)$ is a Banach space with the supremum  norm.

Thus, this is what we have proved so far:

\begin{theorem}
    The disc algebra $\calA =\calC \left( \overline{ \D } \right) \cap H \left( \D \right)$ is a Banach space under the $\sup$ norm
    \begin{align*}
	\norm{f}_{\infty} = \sup_{\abs{z}\le 1} \abs{f\left( z \right)}
    \end{align*}
    \label{thm:disc-algebra-is-B-space}
\end{theorem}

We make a couple of observations at this point:
\begin{enumerate}
    \item Each $f\in \mathcal A$ is the Poisson integral of its boundary values:
\begin{align*}
    f\left( re^{i\theta} \right) = \frac{1}{2\pi} \int_{-\pi}^{\pi} f\left( e^{it} \right) P_{r} \left( e^{i \left( \theta -t \right)} \right) dt
\end{align*}
\item It follows from the Maximum Modulus Theorem that 
    \begin{align*}
	\norm{f}_{\infty} = \sup \abs{f\left( e^{it} \right)}
    \end{align*}
\end{enumerate}

\begin{theorem}[Correspondence of $\calA$ with a closed subspace of $\calC \left( \T \right)$]
    Consider the subspace
    \begin{align*}
	\tilde{\calA} = \left\{ f\in \calC \left( \T \right) \, : \, \int_{-\pi}^{\pi} f\left( e^{it} \right) e^{int}=0  \text{ for } n=1,2,\ldots \right\}
    \end{align*}
    of $\calC \left( \T \right)$. Then there is an isometric isomorphism of $\calA$ with $\tilde{\calA}$.
    \label{thm:correspondence-of-disc-algebra}
\end{theorem}
\begin{proof}
    First, we show that $\tilde{\calA}$ is a closed subspace of $\calC \left( \T \right)$. Let $\left\{ f_{n} \right\}$ be a sequence of functions in $\tilde{\calA}$ converging to $f\in \calC \left( \T \right)$. Consider the following:
    \begin{align*}
	\abs{\int_{-\pi}^{\pi} f\left( e^{it} \right) e^{ikt} dt} &= \abs{\int_{-\pi}^{\pi} f\left( e^{it} \right) e^{ikt} dt - \int_{-\pi}^{\pi} f_n\left( e^{it} \right) e^{ikt} dt}  \\
	&= \int_{-\pi}^{\pi} \abs{f \left( e^{it} \right) -f_{n} \left( e^{it} \right)} dt \\
	& \le 2\pi \norm{f_{n} - f}_{\infty} \to 0 \text{ as } n\to \infty
    \end{align*}
    This shows that $\tilde{\calA}$ is closed under $\calC \left( \T \right)$ with supremum norm.

    Now consider the linear map $T: \calA \to \tilde{\calA}$ given by
    \begin{align*}
	f \stackrel{T}{\longmapsto} f\mid_{\T}
    \end{align*}
    For the sake of convenience, we will write $f\mid_{\T}$ as $f_{\T}$.
    We first need to show this map is well defined! That is, we need to show that
    \begin{equation*}
	\int_{-\pi}^{\pi} f_{\T} \left( e^{it} \right) e^{ikt} dt = 0	
    \end{equation*}
    for all $k\in \N$ but this immediately follows from Cauchy's theorem.

    Note that injectivity is clear from Theorem \ref{thm:Poisson-Integral-Formula}. To show surjectivity, let $f\in \tilde{A}$. We need to show that there is a function $u \in \calA$ such that $u_{\T} = f$. Consider the function
    \begin{equation*}
	u\left( re^{i\theta} \right) =
	\begin{cases}
	    (P*f) (re^{i\theta}) & \text{ if } 0\le r <1 \\ 
	    f\left( e^{i\theta} \right) & \text{ if } r=1
	\end{cases}
    \end{equation*}
    This is the Dirichlet problem on the unit disc! So, $u$ is continuous on $\overline {\D}$. It remains to show that $u$ is analytic on $\D$. But note that for $r\in [0,1)$,
    \begin{align*}
    u\left( re^{i\theta} \right) &= \sum_{n=-\infty}^{\infty} r^{\abs{n}} \hat{f}\left( n \right) e^{int} \\
&= \sum_{n=0}^{\infty} r^{\abs{n}} \hat{f}\left( n \right) e^{int} 
    \end{align*}
    This completes the proof of the theorem!
\end{proof}

In view of the previous theorem, we will simply write $\tilde{\calA}$ as $\calA$.

\begin{definition}
    An analytic trignometric polynomial $p$ on the circle $\T$ is of the form
    \begin{equation*}
	p\left( e^{it} \right) = \sum_{k=0}^{n} a_{k} e^{ik\theta}
    \end{equation*}
    \label{def:trignometric-polynomial}
\end{definition}

\begin{proposition}
    The set of the trignometric polynomials is a dense subset of $\calA$.
    \label{prop:disc-algebra-separable}
\end{proposition}
\begin{proof}
    It is clear that any trignometric polynomial on the circle is a member of the disc algebra. Now, if $f: \T \to \C$ is in $\calA$, then its negative Fourier coefficients are zero! Since, the Cesaro sum of $f$ 
    \begin{equation*}
	s_{n} \left( x \right) = \sum_{k=-n}^{k=n} \hat{f} (n) e^{ikx} = \sum_{k=0}^{n} \hat{f} (n) e^{ikx}
    \end{equation*}
    converge to $f$ uniformly and is a sequence of trignometric polynomial, we are done!
\end{proof}

The following result is used in the proof of the next theorem, so, we prove it here:

\begin{theorem}
    The real parts of functions in $\calA$ are uniformly dense in $\calC \left( \T , \R \right)$. In other words, if $\mu$ is finite signed Borel measure on $\T$ such that $\int f \, d\mu =0$ for every $f\in \calA$ then $\mu$ is the zero measure.
    \label{thm:kinda-Fejer}
\end{theorem}
\begin{proof}
    We first show that any trignometric polynomial of the form
    \begin{equation}
	p\left( e^{it} \right) = \sum_{k=-n}^{n} c_{k} e^{ikt}
	\label{eqn:trig-poly-1}
    \end{equation}
    where $c_{-k} = \overline{c_{k}}$ for each $k\in \left\{ 1,\ldots , n \right\}$ is a real part of a function $f\in \calA$. Note that $p\left( e^{it} \right)$ in Equation \ref{eqn:trig-poly-1} is the real part of the function:
    \begin{equation*}
	f\left( e^{it} \right) = c_{0} + 2c_{1} e^{it} + 2c_{2} e^{2it} + \ldots + 2 c_{n} e^{int}
    \end{equation*}

    
    Now, we claim that every function $f \in C\left( \T , \R \right)$ is a uniform limit of a trignometric polynomial of the form \ref{eqn:trig-poly-1}. We will be done if we show that the negative Fourier coefficients of real valued function is the conjugate of the its positive counterpart, that is, for each $n\in \Z_{\ge 0}$, we have that $\hat{f}\left( -n \right) = \overline{\hat{f} \left( n \right)}$. To show this, take any $n\in \Z _{\ge 0}$ and then observe that
    \begin{align*}
	\hat{f} \left( -n \right) &= \int_{\T} f\left( e^{it} \right) e^{int} \frac{dt}{2\pi} \\
	&= \overline{\int_{\T} f\left( e^{it} \right) e^{-int} \frac{dt}{2\pi}} \\
&= \overline{\hat{f} \left( n \right)}
    \end{align*}

This shows that the Cesaro means of a real valued function is a trignometric polynomial of the form \ref{eqn:trig-poly-1} and since the Cesaro means converges to $f$ uniformly. Thus, the closure of the real parts of $\calA$ is indeed $\calC \left( \T , \R \right)$.

Now, let $\mu$ be a finite signed measure on $\T$ such that $\int f \, d\mu =0$ for every $f \in \calA$. We show that $\mu$ is the zero measure. Notice that if $f\in \calA$ then
\begin{equation*}
    0=\int_{\T} f \, d\mu = \int_{\T} \Re (f) \, d\mu + i\int_{\T} \Im (f) \, d\mu
\end{equation*}
Hence, it follows that 
\begin{equation*}
    \int_{\T} \Re (f) \, d\mu = 0
\end{equation*}
for every $f\in \calA$.
Now, if $g\in \calC \left( \T, \R \right)$ then by the first part of this theorem, there is a sequence $\left\{ f_{n} \right\} \in \calA$ such that $\Re \left( f_{n} \right) $ converges to $g$ uniformly. By the Dominated Convergence Theorem (which holds, thanks to Jordan Decomposition Theorem), we have that $\int_{\T} g \, d\mu = 0$. 

Now to prove that every $\mu = 0$, it suffices to show that $\hat{\mu} \left( n \right) = 0$ for every $n\in \Z$. For a proof of this fact, we refer the reader to take a look at \cite{mashreghi2009representation} or \cite{katznelson2004introduction}. Now, notice that for any $n\in \Z$, we have
\begin{align*}
    \hat{\mu} (n) &= \int_{\T} e^{-int} d\mu \left( e^{it} \right) \\
    &=  \int_{\T} \left(\cos \left( nt \right) - i \sin \left( nt \right) \right) d\mu \left( e^{it} \right) \\
    &= \int_{\T} \cos \left( nt \right) d\mu \left( e^{it} \right) - i\int_{\T} \sin \left( nt \right) d\mu \left( e^{it} \right) \\
    &= 0
\end{align*}
This completes the proof!
\end{proof}

\begin{corollary}
    Let $\mu$ be a finite signed  measure on $\T$ such that $\int f d\mu = 0$ for every $f\in \calA$ which vanishes at the origin then $\mu$ is a constant multiple of Lebesgue measure.
    \label{cor:constant-multiple-of-Lebesgue}
\end{corollary}
\begin{proof}
    We first prove the following claim: If $f \in \calA$ then $\int_{\T} f d\mu = \frac{f\left( 0 \right)}{2\pi}$. Since the negative Fourier cofficients are zero and $f$ is continuous, we have that the Cesaro means converge uniformly to $f$, that is,
    \begin{align*}
	\sum_{k=0}^{\infty} \hat{f} \left( n \right) e^{int} \rightarrow f \text{ uniformly}
    \end{align*}
    Thus,
\begin{align*}
    \int_{T} f\left( e^{it} \right) \frac{dt}{2\pi} &= \frac{1}{2\pi} \int_{\T}  \left(  \sum_{k=0}^{\infty}\hat{f} \left( n \right) e^{int} \right) dt\\
    &= \frac{1}{2\pi} \sum_{k=0}^{\infty} \int_{\T} \left( \hat{f} \left( n \right) e^{int} dt \right) \\
    &= \frac{\hat{f}(0)}{2\pi} 
    \end{align*}

    Now, we proceed to the proof. We define a measure $d\nu = d\mu - \frac{1}{2\pi} \mu \left( \T \right) dt$. Now, we have that
    \begin{align*}
	\int_{\T} f\left( e^{it} \right) d\nu \left( e^{it} \right) &= \int_{\T}[f-f\left( 0 \right)] \left( e^{it} \right) d\mu \left( e^{it} \right) + f\left( 0 \right) \int_{\T} d\nu \left( e^{it} \right) \\
	&= 0
    \end{align*}
    Hence, we have that $d\mu = \frac{1}{2\pi} \mu \left( \T \right) dt$.
\end{proof}

We will be working entirely on the $\T$. So, $\calA$ and $H^{2}$ will be the spaces on the unit circle rather on the open unit disc. 

Now, consider $\calA$ as a subset of $L^{2} \left( \T , \scrB \left( \T \right) , \mu \right)$ where $\mu$ is any finite positive measure. Let $\calA _{0} = \left\{ f \in \calA : \int_{\T} f\left( e^{it} \right) \frac{dt}{2\pi} = \frac{\hat{f} (0)}{2\pi}=0 \right\}$. It is easily seen that $\calA _{0}$ is a subspace of $L^{2} \left( d\mu \right)$. Therefore, we have the closed subspace spanned by $\calA$ is $\llbracket \calA_{0} \rrbracket = \overline{ \operatorname{span} (\calA_{0}) } = \overline{\calA_{0}}$. By a theorem of Hilbert spaces, we have that there is some vector $F \in \llbracket \calA _{0} \rrbracket$ such that
\begin{equation*}
    \inf_{f\in \llbracket \calA _{0} \rrbracket} \int \abs{1-f^{2}} d\mu = \int \abs{1-F^{2}} d\mu
\end{equation*}

But since $d\left( 1, \llbracket \calA_{0}\rrbracket \right) = d \left( 1, \overline \calA_{0} \right) \stackrel{\footnote{this is true for any set in a metric space}}{=} d \left( 1, \calA_{0} \right)$, we have

\begin{equation*}
    \inf_{f\in \calA _{0}} \int \abs{1-f^{2}} d\mu = \int \abs{1-F^{2}} d\mu
\end{equation*}

Note that this $F$ is the orthogonal projection of $1$ into the closed subspace spanned by $\calA_{0}$.

\begin{theorem}
    Let $\mu$ be a finite positive Borel measure on $\T$ and suppose that the constant function $1$ is not in $\llbracket \calA_{0} \rrbracket$. Then let $f=P_{\llbracket \calA_{0} \rrbracket} \left( 1 \right)$. Then the following holds:
    \begin{enumerate}
	\item The measure $d\nu = \abs{1-F^{2}} d\mu$ is a nonzero constant multiple of the Lebesgue measure. In particular, Lebesgue measure is absolutely continuous with respect to $\mu$.
	\item The function $\left( 1-F \right)^{-1} \in H^{2}$.
	\item If $h= \left( \frac{d\mu}{d\theta} \right)$ then $\left( 1-F \right)h \in L^{2} = L^{2} \left( \frac{d\theta}{2\pi} \right)$.
    \end{enumerate}
    \label{thm:prelude-to-fm-riesz}
\end{theorem}

\begin{proof}
    \begin{enumerate}
	\item Let $S=\llbracket \calA _{0} \rrbracket$. We begin to prove part one of the theorem. Let $F= P_{S} \left( 1 \right)$. Then we have by the uniqueness of the decomposition that
    \begin{equation*}
	1=\underbrace{F}_{P_S (1)} + \underbrace{1-F}_{P_{S^{\perp} (1)}}
    \end{equation*}

    Thus, we have that $\left( 1-F \right)$ is orthogonal to every element in $S$ and hence, in particular, any element in $\calA _{0}$ (because $\calA_{0} \subset S \leadsto S^{\perp} \subset \calA_{0} ^{\perp}$). We claim that $1-F$ is orthogonal to $\left( 1-F \right) f $ for every $f\in \calA _{0}$.  But before, we do this, we need to show that $(1-F)f \in L^{2} (d\mu)$. Observe that\begin{align*}
	\int_{\T} \abs{(1-F)f^{2}} d\mu \le \norm{f}_{\infty}^{2} \norm{1-F}_{2}^{2} < \infty
    \end{align*}

    To prove this, note that we showed that $S=\overline{\calA_{0}}$ in the paragraph before the statement of this theoremand since $F \in S$, there is a sequence $\left\{ f_{n} \right\} \in \calA _{0}$ converging to $F$. Hence, we have that $\left\{ f\left( 1-f_{n} \right) \right\}$ is a sequence in $\calA _{0}$\footnote{$\calA_{0}$ is an algebra!} converges to $f\left( 1-F \right)$ in the $L^{2}$-norm. Hence, we have that

    \begin{align*}
	\ip{f(1-F), (1-F)} &= \ip{\lim_{n\to \infty} f\left( 1-f_{n} \right), \left( 1-F \right)} \\
	&= \lim_{n\to \infty} \ip{f\left( 1-f_{n} \right), \left( 1-F \right)} &\text{continuity of the inner product} \\
	&= 0  & \text{$1-F$ is orthogonal to $\calA_{0}$}
    \end{align*}
     
    Now, let $d\nu = \abs{1-F}^{2} d\mu$. We have shown that for any $f\in \calA_{0}$, 
    \begin{align*}
	\int_{\T} f d\nu &= \int_{\T} f \abs{1-F}^{2} d\mu = \ip{f(1-F), (1-F)} = 0
    \end{align*}
    Hence, by Corollary \ref{cor:constant-multiple-of-Lebesgue}, we have that $d\mu = k\, d\lambda$ for some $k \ge 0$.

    Now, we claim that this $k\ne 0$. If $k=0$ then we would have that 
    \begin{equation*}
	\int_{\T} d\nu = 0 \leadsto \int_{\T} \abs{1-F}^{2} d\mu = 0
    \end{equation*}
    Hence, we have that $F=1$ $\mu$-almost everywhere\footnote{See Corollary 2.3.12 of \cite{cohn2013measure}.}. But then we have that $1\in S$ which contradicts our assumption. Hence $k \ne 0$.
\item Observe that part 1 of the theorem tells us that 
    \begin{equation*}
	\abs{1-F}^{2} d\mu = k\, d\lambda \text{ where } k \ne 0
    \end{equation*}
    Then we have that by Lebesgue Decomposition Theorem
    \begin{align*}
	d\mu = h d\lambda + d\mu_{s}
    \end{align*}
    for some positive $\calL^{1}$-function $h$ and some singular measure $\mu_{s}$.
    Hence, we have that
    \begin{align}
	\abs{1-F}^{2} d\mu &= \abs{1-F}^{2} \left( h d\lambda + d\mu_{s} \right) \\
	&= \underbrace{\abs{1-F}^{2} h d\lambda}_{(1)} + \underbrace{\abs{1-F}^{2} d\mu_{s}}_{(2)}
	\label{eqn:decomposed-measure}
    \end{align}
    By the uniqueness of the Lebesgue Decomposition Theorem (one needs to verify that the measures obtained in (1) and (2) are absolutely continuous and singular), we have that
\begin{align*}
    \abs{1-F}^{2} h &= k \text{ $\lambda$-almost everywhere.} \\
    \leadsto \frac{1}{\abs{1-F}^{2}} &= \frac{h}{k} \text{ $\lambda$-almost everywhere.}
\end{align*}
This tells us that $(1-F)^{-1} \in L^{2} \left( d\lambda \right)$ where $\lambda$ is the Lebesgue measure on $\T$. Also, note that $F=1$ $d\mu_{s}$-almost everwhere. Hence, we have that 
\begin{equation}
    \abs{1-F}^{2} d\mu = \abs{1-F}^{2} h d\lambda = k d\lambda \leadsto d\mu = \abs{1-F}^{-2} k d\lambda
    \label{eqn:prelude-to-f-m}
    \end{equation}

    Now, we proceed to show that the negative Fourier coefficients of $\left( 1-F \right)^{-1}$ is 0 then we would have shown that $\left( 1-F \right)^{-1} \in H^{2}$. Consider the following for $f\in \calA_{0}$:
\begin{align*}
    k \int_{\T} \left( 1-F \right)^{-1} f d\lambda &= k \int_{\T} \left( 1-\overline F \right) f \abs{1-F}^{-2} d\lambda \\
    &= k \int_{\T} \left( 1-\overline F \right) f d\mu \\
    &= 0
\end{align*}
The last line follows from the fact that $1-F \perp f$ for any $f\in \calA_{0}$. Hence this holds for any $e^{in\theta}$ and hence we are done.

\item Note that the derivative of $\mu$ with respect to normalized Lebesgue meaures is $h$ as demonstrated in Equation \ref{eqn:decomposed-measure}. Now, Equation \ref{eqn:prelude-to-f-m} shows that 
    \begin{align*}
	\abs{1-F}h = k\abs{1-F}^{-1} \text{ $\lambda$-almost everywhere.}
    \end{align*}
    Since $\abs{ 1-F}^{-1} \in L^{2} \left( \lambda \right)$, so is $(1-F)h$.
    \end{enumerate}
\end{proof}




\begin{corollary}
    Let $\mu$ be a finite positive measure on $\T$ with absolutely continuous part $\mu_{a}$, then
    \begin{align*}
	\inf_{f\in \calA_{0}} \int_{\T} \abs{1-f}^{2} d\mu = \inf_{f\in \calA_{0}} \int_{\T} \abs{1-f}^{2}  d\mu_{a}
    \end{align*}
    where $\mu_{a}$ is the absolutely continuous part of $\mu$.
    In particular, for any singular measure $\mu$, the function $1$ is in the $L^{2} \left( d\mu \right)$ closure of $\calA_{0}$.
\end{corollary}
\begin{proof}
    Let $F$ be the orthogonal projection of $1$ into the closed subspace spanned by $\calA_{0}$. Then we have that 
    \begin{align*}
	\inf_{f\in \calA_{0}} \int_{\T} \abs{1-f}^{2} d\mu = \norm{1-F}^{2} = \int_{\T} \abs{1-F}^2 d\mu
    \end{align*}
    From Equation \ref{eqn:prelude-to-f-m}, we see that
    \begin{equation*}
	\abs{1-F}^{2} d\mu = \abs{1-F}^{2} h d\lambda = \abs{1-F}^{2} d\mu_{a}
    \end{equation*}
    where $h$ is the Radon Nikodym derivative of $\mu$.
    Integrating both sides, we have that
\begin{align*}
	\inf_{f\in \calA_{0}} \int_{\T} \abs{1-f}^{2} d\mu = \inf_{f\in \calA_{0}} \int_{\T} \abs{1-f}^{2}  d\mu_{a}
    \end{align*}
    which is what we wanted to prove.
    Now if $\mu$ is a singular measure then we have that the absolutely continuous part of $\mu$ is the zero measure and hence we have that
\begin{align*}
	\inf_{f\in \calA_{0}} \int_{\T} \abs{1-f}^{2} d\mu = \norm{1-F}^{2} = 0   
    \end{align*}

    From the fact that $\norm{1-F}=0$, we conclude that $1=F$ almost everywhere and hence $1$ is in the closed subspace spanned by $\calA_{0}$ which is in fact $\overline {\calA_{0}}$. This completes the proof of the Corollary.
\end{proof}

\begin{corollary}
    Let $\mu$ be a finite complex Borel measure on $\T$ which is orthogonal to $\calA_{0}$, that is, $\int_{\T} f d\mu =0$ for all $f\in \calA_{0}$. Then the absolutely continuous part and singular parts of $\mu$ are separately orthogonal to $\calA_{0}$. That is, if $\mu_{a}$ and $\mu_{s}$ are the absolutely continuous and the singular parts of the measure $\mu$ then we have that
    \begin{align*}
	\int_{\T} f d\mu_{a} = 0 \text{ and } \int_{\T} f d\mu_{s} = 0
    \end{align*}
    for every $f\in \calA_{0}$. 
\end{corollary}
\begin{proof}
   Let $\rho$ be any finite positive measure satisfying
   \begin{enumerate}
       \item $\mu \ll \rho$ and $\frac{d\mu}{d\rho}$ is bounded, and
       \item $\frac{d\rho}{d\theta} \ge \frac{1}{2\pi}$, that is, the Radon Nikodym derivative with respect to the normalised Lebesgue measure is bounded below by $1/2\pi$.
   \end{enumerate}

   Let $d\mu = \frac{1}{2\pi} h d\theta + d\mu_{s}$ be the Lebesgue decomposition of the measure $\mu$. We define a new measure on $\T$ by 
   \begin{equation*}
       d\rho =  \left( 1+ \abs{h} \right) \frac{d\theta}{2\pi} + d\abs{\mu_{s}}
   \end{equation*}
   where $\abs{\mu_{s}}$ is the variation measure of the complex measure $\mu_{s}$, the singular part of $\mu$.

   We show that $\rho$ is an finite positive measure satisfying properties as in items 1 and 2 above. First, we show that $\rho$ is a finite positive measure. But before that note that the variation measure of any complex measure is  a finite measure. For a proof of the aforementioned fact, we advice the reader to see \cite{cohn2013measure}. Now, consider the following:
   \begin{align*}
       \rho\left( \T \right) &= \int_{\T} d\rho \\
       &= \int_{\T} \left( 1+ \abs{h} \right) \frac{d\theta}{2\pi} + \int_{\T} d\abs{\mu_{s}}  & \text{by definition} \\
       &= 1 + \abs{\mu_{a}} \left( \T \right) + \abs{\mu_{s}} \left( \T \right) & \mu_{a} \text{ is the absolutely continuous part of } \mu \\
       &< \infty
   \end{align*}
   Now we proceed to show that item 1 holds for $\rho$. First, we prove that $\mu \ll \rho$. Let $A \in \scrB \left( \T \right)$ with $\rho \left( A \right) = 0$. Then we have that
   \begin{align*}
       \int_{A} \left( 1+ \abs{h} \right) \frac{d\theta}{2\pi} &=0 \text{ and} \\
       \int_{A} d\abs{\mu_{s}} &= 0
   \end{align*}
   From this, we have that
   \begin{align*}
       \abs{\mu_{a}} \left( A \right) &=  \int_{A} \abs{h} \frac{d\theta}{2\pi} \\
       &\le \int_{A} \left( 1+\abs{h}  \right)\frac{d\theta}{2\pi} \\
       &= 0
   \end{align*}
   Since $\abs{\mu (A)} \le \abs{\mu_{a}} (A)$, we have that $\mu_{a} \left( A \right) = 0$. Also, it follows by the definition of the variation measure that $\mu_s \left( A \right) = 0$. Hence, we have that $\mu $ is absolutely continuous with respect to $\rho$. Now, we need to show that $\frac{d\mu}{d\rho}$ is bounded.


We proceed to show that item 2 holds. Observe that
\begin{align*}
    \frac{d\rho}{d\theta} = \left( 1+\abs{h} \right) \frac{1}{2\pi} \ge \frac{1}{2\pi}
\end{align*}
and hence we are done showing that $\rho$ is such a measure.

Now, let $f\in \calA_{0}$. Then we have that
\begin{align*}
    \int_{\T} \abs{1-f}^{2} d\rho &= \int_{\T} \abs{1-f}^{2} \frac{d\rho}{d\theta} d\theta + \int_{\T} \abs{1-f}^2 d\mu_{s} \\
    &\ge \int_{\T} \abs{1-f}^{2} \frac{d\rho}{d\theta} d\theta \\
    & \ge \int_{\T} \abs{1-f}^{2} \frac{d\theta}{2\pi} \\ 
    &= \norm{1-f}^{2} \\
    &= \norm{1}^{2} + \norm{f}^{2} & \ip{f,1} = 0  \text{ as } f\in \calA_{0} \\
&\ge 1 
\end{align*}

Now, let $F$ be the orthogonal projection of $1$ into the closed subspace spanned by $\calA_{0}$ of $L^{2} \left( d\rho \right)$. Thus, by definition, we have that 
\begin{equation*}
    \int_{\T} \abs{1-F}^{2} = \inf_{f\in \calA_0} \int_{\T} \abs{1-f}^{2} d\rho \ge 1
\end{equation*}
Hence, we have that $1$ is not in the closed subspace spanned by $\calA_{0}$ of $L^{2} (d\rho)$. Hence, we have by Theorem \ref{thm:prelude-to-fm-riesz} that $\left( 1-F \right)^{-1} \in H^{2}$ and that 
\begin{equation*}
    \left( 1-F \right) \left( 1+\abs{h} \right) \in L^{2} = L^{2} \left( \frac{d\theta}{2\pi} \right)
\end{equation*}
From this, we can conclude that $h \in L^{2}$ by the following fashion:
\begin{equation*}
    \abs{h} = \underbrace{\left( 1-F \right) \left( 1+ \abs{h} \right)}_{\in L^{2}} \underbrace{\left( \abs{1-F}^{-1} \right)}_{\in H^{2} \leadsto \in L^{2}} -1
\end{equation*}
Hence, we have that $(1-F)h \in L^{2}$.

Suppose $g\in \calA_{0}$. We claim that
\begin{align*}
    \int_{\T} \left( 1-F \right) g d\mu = 0
\end{align*}
\end{proof}

Since $F\in \overline \calA_{0}$, there is a sequence $f_{n} \in \calA_{0}$ such that $f_{n} \to F$ in $L^{2} (d\rho)$. Hence,
\begin{align*}
    (1-f_{n})g \to \left( 1-F \right)g
\end{align*}
in $L^{2} \left( d\rho \right)$.

Since $\frac{d\mu}{d\rho}$ is bounded, we have that 
\begin{equation*}
    \int_{\T} \left( 1-F \right) g d\mu = \int_{\T} \left( 1-F \right) g \frac{d\mu}{d\rho} d\rho = \lim_{n\to \infty} \int_{\T} \left( 1-f_{n} \right) \frac{d\mu}{d\rho} d\rho = 0 \text{.}
\end{equation*}

Now, $1-F = 0$ with respect to $d\abs{\mu_{s}}$ which implies that $1-F= 0$ with respect to $d\mu_{s}$.

Thus, we have that 
\begin{align*}
    \left( 1-F \right) d \mu = \frac{1}{2\pi} \left( 1-F \right)h d\theta
\end{align*}

Hence, we have for $g \in \calA_{0}$
\begin{align*}
    \int_{\T} \left( 1-F \right) gh d\theta &= \int_{\T} (1-F)g d\mu  \\
    &= 0
\end{align*}

Since $\overline{A} = H^{2}$. There is a sequence of functions ${g_{n}}$ in $\calA$ converging to $\left( 1-F \right)^{-1} \in L^{2} \left( \frac{d\theta}{2\pi} \right)$ in $L^{2}$-norm. Hence, we have for any $f\in \calA_{0}$,
\begin{equation*}
    \int_{\T} g_{n}f\left( 1-F \right)h d\theta = 0
\end{equation*}


Also, note that $\left( 1-F \right)h \in L^{2}$. From this, we can conclude that 
\begin{equation*}
    \int_{\T} fh d\theta = 0
\end{equation*}

\begin{theorem}[Frigyez and Marcel Riesz, 1916]
    Let $\mu$ be a complex Borel measure on the unit circle $\T$ such that
    \begin{equation*}
	\int_{\T} e^{in\theta} d\mu \left( e^{i\theta} \right) =0
    \end{equation*}
    for all $n\in\N$. Then $\mu$ is absolutely continuous with respect to the Lebesgue measure.
    \label{thm:f-and-m-riesz}
\end{theorem}
\begin{proof}
    Suppose that $\mu$ is orthogonal to $\calA_{0}$. Then from the previous corollary, we have that $\mu_{a}$ is orthogonal to $\calA_{0}$ and $\mu_{s}$ is orthogonal to $\calA_{0}$. Since $\abs{\mu_{s}}$ is singular (definition?), we can find a sequence of functions $\left\{ f_{n} \right\}$ converging to $1$ in $L^{2} \left( d\abs{\mu_{s}} \right)$. Since $\mu_{s}$ is orthogonal to $\calA_{0}$, we have that
    \begin{equation*}
	\int_{\T} d\mu_{s} = \lim_{n\to \infty} \int_{\T} f_{n} d\mu_{s} = 0
    \end{equation*}
    Thus, $\mu_{s}$ is orthogonal to $1$. The singular measure $e^{-i\theta} d\mu_{s}$ is now orthogonal to $\calA_{0}$; hence it is orthogonal to $1$, that is, 
    \begin{equation*}
    \int_{\T} e^{-i\theta} d\mu_{s} = 0
    \end{equation*}
    Similarly, $e^{-2i\theta} d\mu_{s}$ is orthogonal to $\calA_{0}$, and consequently, orthogonal to $1$. Repeating thei process, we conclude that
    \begin{equation*}
	\int_{\T} e^{in\theta} = 0
    \end{equation*}
    for $n\in \Z$.
    Hence $\mu_{s}$ must be the zero measure. Therefore our original measure $\mu$ is absolutely continuous. 
\end{proof}
\section{Szegö's Theorem}
\begin{theorem}[Szegö]
    Let $h \in L^{1} \left( \T \right)$ and suppose that $h \ge 0$ on $\T$ (that makes $h$ real valued!). Then the following holds:
    \begin{align*}
	\exp \left[ \frac{1}{2\pi} \int_{-\pi}^{\pi} \log h\left( e^{it} \right) dt \right] = \inf_{f\in \calA_{0}} \frac{1}{2\pi} \int_{-\pi}^{\pi} he^{\Re \left( f \right)}dt
    \end{align*}
    \label{thm:prelude-to-szego-kolmorogoff-krein}
\end{theorem}
\begin{proof}
    We first claim that $\log h$ is integrable iff $\int \log h d\theta = -\infty$. To see this, first we decompose 
    \begin{align*}
	\int \log h d\theta = \underbrace{\int \left( \log h \right)^{+} d\theta}_{(\romannumeral1)} - \underbrace{\int \left( \log h \right)^{-} d\theta}_{(\romannumeral2)}
    \end{align*}
First observe that $\log h \le h$ (this follows by definition of $\log$) and hence, we have that $\left( \log h \right)^{+} \le h^{+}=h$. Thus, we have that integral $\left( \romannumeral1 \right)$ is finite. Therefore, $\log h$ is not integrable iff $\int \log h d\theta =- \infty$. 

    Regardless of integrability of $\log h$, we have the following relation:
    \begin{equation}
	\exp \left[ \frac{1}{2\pi} \int_{\T} \log h\left( e^{it} \right) dt \right] \le \frac{1}{2\pi} \int_{\T} h\left( e^{it} \right) dt
	\label{eqn:almost-am-gm}
    \end{equation}

    One can prove the famililar arithmetic and geometric means inequality from the above inequality. Let $n\in \N$ be arbitrary and $a_{1}, \ldots, a_{n}$ be positive real numbers. Let $A_{1}, A_{2}, \ldots , A_{n}$ be the disjoint subsets of $\T$ which cover it. Therefore, we have that
    \begin{align*}
	\mu \left( A_{i} \right) = \frac{2\pi}{n} \text{ for all } i=1,2,\ldots,n\text.
    \end{align*}
    Observe the following:
    \begin{align*}
	\log\left( a_{1} \chi_{A_{1}} + \ldots + a_{n} \chi_{A_{n}} \right) &= \log \left( a_{1} \chi_{A_{1}} + \ldots + a_{n} \chi_{A_{n}} \right) \\
	&= \log \left( a_{1} \right) \chi_{A_{1}} + \ldots + \log \left( a_{n} \right) \chi_{A_{n}}
    \end{align*}
    Integrating both sides, we get
    \begin{align*}
	\frac{1}{2\pi} \int_{\T} \log\left( a_{1} \chi_{A_{1}} + \ldots + a_{n} \chi_{A_{n}} \right) dt &= \frac{1}{2\pi} \int_{\T}  \log \left( a_{1} \right) \chi_{A_{1}} + \ldots + \log \left( a_{n} \right) \chi_{A_{n}} dt \\
	&= \frac{1}{2\pi} \left( \log \left( a_{1} \right) \mu \left( A_{1} \right) + \ldots +  \log \left( a_{n} \right) \mu \left( A_{n} \right) \right) \\
	&=  \log \left( \sqrt[n]{a_{1}\ldots a_{n}} \right)
    \end{align*}
     
    It can be seen now that it follows immediately.


    Now, if $h\in L^{\T}$ and is nonnegative, $g$ is any real valued continuous function on $\T$ then we have that $he^{g}$ is an nonnegative $L^{1} \left( \T \right)$ because product of an $L^{1}$ function with a bounded funtion on a finite measure space retains it back and therefore, we can apply Equation \ref{eqn:almost-am-gm}. Futhermore, if we assume that $\int_{\T} g dt =0 $ then we have that 
    \begin{align*}
	\exp\left[ \frac{1}{2\pi} \int_{\T} \log \left( h\left( e^{it} \right) e^{g\left( e^{it} \right)}\right) dt \right] \le \frac{1}{2\pi} \int_{\T} h\left( e^{it} \right) e^{g\left( e^{it} \right)} dt \\
	\leadsto \exp\left[ \frac{1}{2\pi} \int_{\T} \log \left( h\left( e^{it} \right) \right) dt + \cancelto{0}{\int_{\T} g\left( e^{it} \right) dt} \right] \le \frac{1}{2\pi} \int_{\T} h\left( e^{it} \right) e^{g\left( e^{it} \right)} dt \\
\leadsto \exp\left[ \frac{1}{2\pi} \int_{\T} \log \left( h\left( e^{it} \right) \right) dt  \right] \le \frac{1}{2\pi} \int_{\T} h\left( e^{it} \right) e^{g\left( e^{it} \right)} dt \\
    \end{align*}

    Now, let $g=\Re f$ for some $f\in \calA_{0}$. Then we have that $g$ is continuous and real valued and hence the previous inequality holds, therefore, we have that 

    \begin{equation*}
	\exp\left[ \frac{1}{2\pi} \int_{\T} \log \left( h\left( e^{it} \right) \right) dt  \right] \le \frac{1}{2\pi} \int_{\T} h\left( e^{it} \right) e^{\Re \left( f\left( e^{it} \right) \right)} dt \text{.}
    \end{equation*}
    Since $f \in \calA_{0}$ is arbitrary, we have that 
    \begin{equation*}
	\exp\left[ \frac{1}{2\pi} \int_{\T} \log \left( h\left( e^{it} \right) \right) dt  \right] \le \inf_{f\in\calA_{0}} \frac{1}{2\pi} \int_{\T} h\left( e^{it} \right) e^{\Re \left( f\left( e^{it} \right) \right)} dt \text{.}
    \end{equation*}

    Now, we claim the set of functions $\Re f$ where $f\in \calA$ is dense subset of $L^{1}_{\R}\left( \T \right)$. This is easy to see. Let $f\in L^{1}_{\R} \left( \T \right)$. Then denseness of continuous functions in $L^{1} _{\R}\left( \T \right)$, we have that there is some function $g\in \calC _{\R} \left( \T \right)$ such that 
    \begin{align*}
	\norm{f-g}_{1} < \frac{\varepsilon}{2}
    \end{align*}
    Also, by Theorem \ref{thm:kinda-Fejer}, we have that there is some function $h\in \calA$ such that 
    \begin{align*}
	\norm{g-\Re h}_{\infty} < \frac{\varepsilon}{2}
    \end{align*}
    Therefore, we have that
    \begin{align*}
	\norm{f-\Re h}_{1} & \le \norm{f-g}_{1} + \norm{g-\Re h} & \text{Minkowski inequality} \\
	& \le \norm{f-g}_{1} + \norm{f-g}_{\infty} & \text{finite measure space} \\
	& \le \frac{\varepsilon}{2} + \frac{\varepsilon}{2} = \varepsilon\text{.}
    \end{align*}

    This shows the denseness.

    Now, we show that the $\left\{ f\in L^{1}_{\R} \left( \T \right) : \int_{\T} f d\theta = 0 \right\}$ is dense in set of all functions $f \in L^{1}_{\R} \left( \T \right)$ such that $\int f d\theta = 0$. This is immediate because closure distributes over intersections and $\left\{ f\in L^{1}_{\R} \left( \T \right) : \int_{\T} f d\theta = 0 \right\}$ is a closed subspace of $L^{1}_{\R} \left( \R \right)$.

    Now, we claim that $\underbrace{\inf \{ \int he^{g} d\theta : g \in L^{1}, \int g =0 \}}_{:=\alpha} = \underbrace{\inf \left\{ \int he^{g} d\theta : g \in L^{\infty} , \int g =0 \right\}}_{:=\beta}$. Since $L^{\infty} \subset L^{1}$, we have that $\alpha \le \beta$ is clear. To show the reverse inequality, let $g \in L^{1}$ with $\int g d\theta = 0$. By the Simple Function Approximation Theorem, there is a sequence $g_{n} \in L^{\infty}$ such that $0 \le \abs{g_1} \le \abs{g_2} \le \ldots \le \abs{g}$ and $g_{n}$ converging to $g$ pointwise.

\end{proof}

\begin{theorem}[Szegő, Kolmogoroff-Krein]
    Let $\mu$ be a finite Borel measure on the unit circle $\T$ and $h$ be the derivative of $\mu$ with respect to the normalised Lebesgue measure. Then
    \begin{equation*}
	\inf_{f\in \calA_{0}} \inf \abs{1-f}^{2} \, d\mu = \exp  \left[ \frac{1}{2\pi} \int_{\pi}^{\pi} \log h\left( e^{i\theta} \right) \, d\theta \right]
    \end{equation*}
    \label{thm:szego-kolmogoroff-krein}
\end{theorem}
Most of the work to prove this theorem has already been done. We relegate it to the references. For a proof of this, the reader may take a look at \cite{hoffman2007banach}.

\section{Completing the Discussion of \texorpdfstring{$H^1$}{\text{H1}}}

Let $F\in H^{1} \left( \D \right)$. We showed in Chapter $1$ that $F$ is a Poisson integral of a complex measure. That is,
\begin{equation*}
    F=P\left[ \mu \right]
\end{equation*}
for some $\mu \in \calM \left( \T \right)$. Furthermore, we showed that the measures
$d\mu_{r} = \frac{1}{2\pi} F_{r} dt$ converge to $d\mu$ in the weak* topology of measures.

As promised earlier, we have shown that an analytic measure $\mu$ is absolutely continuous that 
\begin{equation*}
    d\mu = \frac{1}{2\pi} \tilde{f} dt
\end{equation*}
for some $f\in L^{1} \left( \T \right)$.

By the previous paragraph, we have that $F_{r}$ converges to $f$ in $L^{1}$ and we have that 
\begin{equation*}
    \tilde{f} \left( e^{i\theta} \right) = \lim_{z\to e^{i\theta}} f\left( z \right)
\end{equation*}
exists for almost all $e^{i\theta} \in \T$. Thus, we may identity this boundary function $\tilde {f} \in H^{1} \left( \T \right)$ with $F\in H^{1} \left( \D \right)$.

The proof of the following theorem relies on Szego's theorem. We state the theorem with proof and the proof of the theorem can be found in \cite{hoffman2007banach}.

\begin{theorem}
    Let $f\in H^{1} \left( \T \right)$ be nonzero then the function $\log \abs{f\left( e^{i\theta} \right)}$ is integrable and we have that 
    \begin{equation*}
	\frac{1}{2\pi} \int_{\T} \log \abs{f} d\theta \ge \log \abs{f\left( e^{i0} \right)}.
    \end{equation*}
    \label{thm:szego-integrable-h1}
\end{theorem}

\begin{corollary}
    If $f \in H^{1} \left( \T \right)$, $f$ cannot vanish on a set of positive Lebesgue measure unless $f \equiv 0$ on $\T$
    \label{cor:H1-f-zero}
\end{corollary}
\begin{proof}
    Assume $f\not\equiv 0$. We wish to show that $f$ cannot vanish on a set of positive Lebesgue measure on $\T$. Since $\log \abs{f}$ is integrable on $\T$ (by Szego's theorem), we have that 
    \begin{equation*}
	\lambda \left( \left\{ e^{i\theta} \in \T \, : \, \abs{\log \abs{f\left( e^{i\theta}  \right)}}  = +\infty\right\} \right) = 0
    \end{equation*}
    since integrability implies finite almost everywhere. Therefore we have that 
    \begin{equation*}
	\lambda \left( \left\{ e^{i\theta} \in \T \, : \, f\left( e^{i\theta} \right) = 0 \right\} \right) = \lambda \left( \left\{ e^{i\theta} \in \T \, : \, \log \abs{f\left( e^{i\theta}  \right)}  = -\infty\right\} \right) = 0
    \end{equation*}
   
    Thus, if $S$ is any measurable subset of the $\T$ where $f$ vanishes then we must be have that $\lambda \left( S \right) = 0$. This shows that $f$ cannot vanish on a set of positive Lebesgue measure.
\end{proof}

\end{document}
